% Options for packages loaded elsewhere
\PassOptionsToPackage{unicode}{hyperref}
\PassOptionsToPackage{hyphens}{url}
\PassOptionsToPackage{dvipsnames,svgnames*,x11names*}{xcolor}
%
\documentclass[
  11pt,
]{article}
\usepackage{lmodern}
\usepackage{amssymb,amsmath}
\usepackage{ifxetex,ifluatex}
\ifnum 0\ifxetex 1\fi\ifluatex 1\fi=0 % if pdftex
  \usepackage[T1]{fontenc}
  \usepackage[utf8]{inputenc}
  \usepackage{textcomp} % provide euro and other symbols
\else % if luatex or xetex
  \usepackage{unicode-math}
  \defaultfontfeatures{Scale=MatchLowercase}
  \defaultfontfeatures[\rmfamily]{Ligatures=TeX,Scale=1}
\fi
% Use upquote if available, for straight quotes in verbatim environments
\IfFileExists{upquote.sty}{\usepackage{upquote}}{}
\IfFileExists{microtype.sty}{% use microtype if available
  \usepackage[]{microtype}
  \UseMicrotypeSet[protrusion]{basicmath} % disable protrusion for tt fonts
}{}
\makeatletter
\@ifundefined{KOMAClassName}{% if non-KOMA class
  \IfFileExists{parskip.sty}{%
    \usepackage{parskip}
  }{% else
    \setlength{\parindent}{0pt}
    \setlength{\parskip}{6pt plus 2pt minus 1pt}}
}{% if KOMA class
  \KOMAoptions{parskip=half}}
\makeatother
\usepackage{xcolor}
\IfFileExists{xurl.sty}{\usepackage{xurl}}{} % add URL line breaks if available
\IfFileExists{bookmark.sty}{\usepackage{bookmark}}{
\usepackage{hyperref}
}
\hypersetup{
  pdftitle={Systèmes d'exploitation},
  pdfauthor={Première NSI, Lycée du Parc},
  colorlinks=true,
  linkcolor=Maroon,
  filecolor=Maroon,
  citecolor=Blue,
  urlcolor=Blue,
  pdfcreator={LaTeX via pandoc}}
\urlstyle{same} % disable monospaced font for URLs
\usepackage[top=20mm,left=20mm,right=20mm,heightrounded]{geometry}
\usepackage{longtable,booktabs}
% Correct order of tables after \paragraph or \subparagraph
\usepackage{etoolbox}
\makeatletter
\patchcmd\longtable{\par}{\if@noskipsec\mbox{}\fi\par}{}{}
\makeatother
% Allow footnotes in longtable head/foot
\IfFileExists{footnotehyper.sty}{\usepackage{footnotehyper}}{\usepackage{footnote}}
\makesavenoteenv{longtable}
\usepackage{graphicx}
\makeatletter
\def\maxwidth{\ifdim\Gin@nat@width>\linewidth\linewidth\else\Gin@nat@width\fi}
\def\maxheight{\ifdim\Gin@nat@height>\textheight\textheight\else\Gin@nat@height\fi}
\makeatother
% Scale images if necessary, so that they will not overflow the page
% margins by default, and it is still possible to overwrite the defaults
% using explicit options in \includegraphics[width, height, ...]{}
\setkeys{Gin}{width=\maxwidth,height=\maxheight,keepaspectratio}
% Set default figure placement to htbp
\makeatletter
\def\fps@figure{htbp}
\makeatother
\setlength{\emergencystretch}{3em} % prevent overfull lines
\providecommand{\tightlist}{%
  \setlength{\itemsep}{0pt}\setlength{\parskip}{0pt}}
\setcounter{secnumdepth}{5}

\title{Systèmes d'exploitation}
\usepackage{etoolbox}
\makeatletter
\providecommand{\subtitle}[1]{% add subtitle to \maketitle
  \apptocmd{\@title}{\par {\large #1 \par}}{}{}
}
\makeatother
\subtitle{Thème architectures matérielles et systèmes d'exploitation}
\author{Première NSI, \href{https://frederic-junier.org/}{Lycée du
Parc}}
\date{}

%%%jolis boites

\usepackage{fancybox, graphicx}



%%%%%%%%%%%%%%%%Packages et Macros Frederic%%%%%%%%%%%%%%%%%%%%%%%%%%%%%


%%%%Insertion de liens hypertextes %%%%

            
%%%%%%%%%%PSTricks%%%%%%%%%%%%

\usepackage{pstricks,pst-plot,pst-text,pst-tree,pst-eps,pst-fill,pst-node,pst-math,pstricks-add,pst-xkey,pst-eucl}


%%%%%%%Tikz%%%%%%%%%%%%%%%
\usepackage{pgf,tikz,tkz-tab}
% Pour les tableaux de signes ou de variations avec tkz-tab voir https://zestedesavoir.com/tutoriels/439/des-tableaux-de-variations-et-de-signes-avec-latex/#1-13389_tikz-un-package-qui-en-a-dans-le-ventre
\usetikzlibrary{arrows}
\usetikzlibrary{shapes.geometric}
\usetikzlibrary{shapes.geometric}
\usetikzlibrary{petri}
\usetikzlibrary{decorations}
\usetikzlibrary{arrows}
\usetikzlibrary{math}
 %Variables must be declared in a tikzmath environment but
       % can be used outside
%       \tikzmath{int \n; \n = 508; \x1 = 1; \y1 =1; 
%                   %computations are also possible
%                    \x2 = \x1 + 1; \y2 =\y1 +3; } 


%%%%%%%%%%%%%%%%%%%%%%%%%%%%%%%%%%%%%%%%
%%%%%%%%%%%Commandes Tikz Perso%%%%%%%%%%%%%%%

% Définition des nouvelles options xmin, xmax, ymin, ymax
% Valeurs par défaut : -3, 3, -3, 3
\tikzset{
xmin/.store in=\xmin, xmin/.default=-3, xmin=-3,
xmax/.store in=\xmax, xmax/.default=3, xmax=3,
ymin/.store in=\ymin, ymin/.default=-3, ymin=-3,
ymax/.store in=\ymax, ymax/.default=3, ymax=3,
}
% Commande qui trace la grille entre (xmin,ymin) et (xmax,ymax)
\newcommand {\grille}[2]
{\draw[help lines,black, thick] (\xmin,\ymin) grid[xstep=#1, ystep=#2] (\xmax,\ymax);}
% Commande \axes
\newcommand {\axes} {
\draw[->,very thick] (\xmin,0) -- (\xmax,0);
\draw[->,very thick] (0,\ymin) -- (0,\ymax);
\draw (0.95*\xmax, 0) node[above] {};
\draw (0, 0.95*\ymax) node[left] {};
}
% Commande qui limite l?affichage à (xmin,ymin) et (xmax,ymax)
\newcommand {\fenetre}
{\clip (\xmin,\ymin) rectangle (\xmax,\ymax);}

%Exemple d'utilisation

%\begin{center}
%\begin{tikzpicture} [xmin=-2,xmax=2,ymin=0,ymax=5]
%\grille{1} \axes \fenetre
%\draw plot[smooth] (\x,\x^2);
%\end{tikzpicture}
%\end{center}

%style pour la perspective cavalière française
%voir Tikz pour l'impatient page 68
\tikzset{math3d/.style=
{x= {(-0.353cm,-0.353cm)}, z={(0cm,1cm)},y={(1cm,0cm)}}}

%%%%%%%Symbole pour code calculatrice%%%%%%

%Flèche remplie pour défilement de menu

\newcommand{\flechefillright}{
\begin{tikzpicture}[scale=0.15] \fill (0,0)--(2,1)--(0,2)--cycle;
\end{tikzpicture}}

%%%%%%%%%%%%%Symboles pour calculatrice Casio%%%%
\newcommand{\execasio}{\Pisymbol{psy}{191}} %Retour chariot
\newcommand{\dispcasio}{\begin{pspicture}(.1,.1)\pspolygon*(.1,0)(.1,.1)\end{pspicture}} %Triangle « Disp »
\newcommand{\dispcasiotikz}{\begin{tikzpicture}[scale=0.2]
\fill (0,0) -- (1,0) -- (1,1) -- cycle;
\end{tikzpicture}} %Triangle « Disp »
%

%Fleche entre deux lignes, d'apres 'un bon petit' : http://forum.mathematex.net/latex-f6/fleches-entre-deux-lignes-pour-resolution-d-equation-t10283.html#p99817
\newcommand\addnode[1]{\Rnode{#1}{}}
\newcommand\linknode[3]{\ncbar[angleA=0,angleB=0,nodesep=1ex,arm=10ex,offset=-2pt]{->}{#1}{#2}\Aput{\vphantom{x}#3}}


%%Commande pour touche de calculatrice

\newcommand\tc[1]{%
{
\begin{tikzpicture}
\node[draw,rectangle,rounded corners=3pt] (P) at (0,0){#1};
\end{tikzpicture}
}
}

%%%%%%%%%%%%%%%%%%%%%%%%%%%%%%%%%%%%%%%%
%%%%%%%%%%%Fin Commandes Tikz%%%%%%%%%%%%%%%


%%%%%%%%%%%%Specifiques%%%%%%%%%%%
\usepackage{wrapfig}
%pour insérer une figure à droite ou à gauche d'un texte
%\begin{wrapfigure}[nb lignes]{placement l,r,c,i(inside),o(outside)}[overhang]{width}
%ce package fonctionne mal à proximité des listes
%%%%%%%%%%%%%%%%%%%%%%%%%%%%%%%%%%%%%

%%%%%Environnements et symboles spéciaux pour faire joli%%%%%%

%%%Bclogo, pour des environnements + jolis avec insertion de logo%%%%
%Dépendances de  bclogo
\usepackage{xkeyval}  
\usepackage{etoolbox}
\usepackage{ifpdf}
\usepackage[framemethod=tikz]{mdframed}
\usepackage[tikz]{bclogo}

%\newcommand\bcpython{\includegraphics[width=17pt]{/home/fjunier/Maths/python-logo.eps}}
\newcommand\bcpython{\includegraphics[width=17pt]{/home/fjunier/Maths/python-logo.png}}
%\newcommand\bcpython{\includegraphics[width=17pt]{/home/frederic/Maths/python-logo.png}}

%% Framed
\usepackage{framed}  %Le package « framed» Crée 3 nouveaux environnements, qui se comportent comme des minipage de largeur \linewidth, mais permettant en plus de se casser entre plusieurs pages.     * framed : avec un cadre autour;     * shaded : avec un fonc coloré (il faut définir la couleur shadecolor);     * leftbar : avec une barre le long du côté gauche.

%%%%%%%%%%%%%%%%%%%Présentation de codes sources%%%%%%%%%%%%%%%%%
\usepackage{listings}
%On utilise l?environnement lstlisting pour insérer
%un code source.
%En plus de l?environnement lstlisting, on peut également utiliser la
%commande \lstinline qui fonctionne comme la commande \verb, en ce
%sens qu?on peut utiliser n?importe quel caractère comme délimiteur. Enfin,
%la commande \lstinputlisting permet de charger un code source depuis
%un fichier externe.
%Il y a deux manières de préciser des options : soit via l?option de l?envi-
%ronnement ou de la commande, soit en utilisant la commande \lstset
%qui permet de définir des options de manière globale.

\lstset{ %
  language=Python,                % the language of the code
  basicstyle=\ttfamily,           % the size of the fonts that are used for the code
  %numbers=left,                   % where to put the line-numbers
  numberstyle=\tiny,  % the style that is used for the line-numbers
  %stepnumber=2,                   % the step between two line-numbers. If it's 1, each line 
                                  % will be numbered
  %numbersep=5pt,                  % how far the line-numbers are from the code
  backgroundcolor=\color{white},      % choose the background color. You must add \usepackage{color}
  showspaces=false,               % show spaces adding particular underscores
  showstringspaces=false,         % underline spaces within strings
  showtabs=false,                 % show tabs within strings adding particular underscores
  frame=single,                   % adds a frame around the code
  rulecolor=\color{black},        % if not set, the frame-color may be changed on line-breaks within not-black text (e.g. comments (green here))
  tabsize=4,                      % sets default tabsize to 2 spaces
  captionpos=b,                   % sets the caption-position to bottom
  breaklines=true,                % sets automatic line breaking
  breakatwhitespace=false,        % sets if automatic breaks should only happen at whitespace
  %title=\lstname,                   % show the filename of files included with \lstinputlisting;
                                  % also try caption instead of title
  breakindent=1cm,
  keywordstyle=\color{blue},          % keyword style
  commentstyle=\color{red},       % comment style
  %stringstyle=\ttfamily\color{green},         % string literal style
  escapeinside={\%*}{*)},            % if you want to add LaTeX within your code
  morekeywords={*,...},              % if you want to add more keywords to the set
  deletekeywords={...}              % if you want to delete keywords from the given language
  upquote=true,columns=flexible,
xleftmargin=1cm,xrightmargin=1cm,
 inputencoding=utf8,			%Les lignes qui suivent sont pour le codage utf8
  extendedchars=true,
  literate=%
            {é}{{\'{e}}}1
            {è}{{\`{e}}}1
            {ê}{{\^{e}}}1
            {ë}{{\¨{e}}}1
            {û}{{\^{u}}}1
            {ù}{{\`{u}}}1
            {â}{{\^{a}}}1
            {à}{{\`{a} }}1
            {î}{{\^{i}}}1
            {ô}{{\^{o}}}1
            {ç}{{\c{c}}}1
            {Ç}{{\c{C}}}1
            {É}{{\'{E}}}1
            {Ê}{{\^{E}}}1
            {À}{{\`{A}}}1
            {Â}{{\^{A}}}1
            {Î}{{\^{I}}}1
}

\lstdefinestyle{rond}{
  numbers=none,
  backgroundcolor=\color{gristclair},
  frameround =tttt
}

\lstdefinestyle{compil}{
  numbers=none,
  backgroundcolor=\color{gristclair}
}
%\lstset{language=Python,basicstyle=\small , frame=single,tabsize=4,showspaces=false,showtabs=false,showstringspaces=false,numbers=left,numberstyle=\tiny , extendedchars=true}



%%%%%%%%%%%%%%%%%%%%%%%%%%%%%%%%%%%%%%%%%%%%%%%%%%%%%%%%%%%%%%%%%%%%%%%%
%%%%%%%%%%%%%%%%%%%%Environnements persos%%%%%%%%%%%%%%%%%%%%%%%%%%%%%%%%
%Syntaxe :
%\newenvironment{nom}[nombre d'args][defaut]{definitions initiales}{definitions finales}
%definitions intiales sont les commandes appelées par \begin{nom}
%Definitions finales sont les commandes appelées par \end{nom}

%%%%%%%%%%%%%%%%Définitions des environnemts persostheoreme, exemple ..%%%%
%%%% Exercice avec encadré %%%%
\newcounter{exo}
\newenvironment{exercice}[1]
{\par \medskip   \addtocounter{exo}{1} \noindent  
\begin{bclogo}[arrondi =0.1,   noborder = true, logo=\bccrayon, marge=4]{~\textbf{Exercice} \textbf{\theexo} {\itshape #1} }  \par}
{
\end{bclogo}
 \par \bigskip }

%%Axiomes, Theoremes, Propriété, Définition, Methode, Preuve


\newenvironment{axiome}[1]
{\par \medskip   \begin{leftbar} \noindent \underline{\textbf{Axiome}}\hspace{0.5cm}{\itshape #1}   \vspace*{10pt} \par }
{\end{leftbar}  \par \medskip }


\newcounter{thme}
\newenvironment{theoreme}[1]
{\par \medskip  \addtocounter{thme}{1} \noindent  
\begin{bclogo}[arrondi =0.1,  ombre = true, barre=none, logo=\bcbook, marge=4]{~\textbf{Théorème} \textbf{\thethme} {\itshape #1} }   \par}
{
\end{bclogo}
 \par \bigskip}

 \newenvironment{theoremedef}[1]
{\par \medskip   \addtocounter{thme}{1} \noindent  
\begin{bclogo}[arrondi =0.1,  ombre = true, barre=none, logo=\bcbook, marge=4]{~\textbf{Théorème-Définition} \textbf{\thethme} {\itshape #1} }   \par}
{
\end{bclogo}
 \par \bigskip }
 
\newcounter{prop}
\newenvironment{propriete}[1]
{\par \medskip   \addtocounter{prop}{1} \noindent  
\begin{bclogo}[arrondi =0.1,  ombre = true, barre=none, logo=\bcbook, marge=4]{~\textbf{Propriété} \textbf{\theprop} {\itshape #1} }   \par}
{
\end{bclogo}
 \par \bigskip }


\newenvironment{corollaire}[1]
{\par \medskip   \noindent  
\begin{bclogo}[arrondi =0.1,  ombre = true, barre=none, logo=\bcbook, marge=4]{~\textbf{Corollaire} {\itshape #1} } \par }
{
\end{bclogo}
 \par \bigskip }

\newenvironment{demo}[1]
{\par \medskip   \noindent  
\begin{bclogo}[arrondi =0.1,  ombre = true, barre=zigzag, noborder = true, logo=\bcloupe, marge=0]{~\textbf{Démonstration} {\itshape #1} } \par \vspace{10pt}}
{
\end{bclogo}
 \par \bigskip }

\newcounter{activite}
\newenvironment{activite}[1]
{\par \medskip   \noindent   \addtocounter{activite}{1}
\begin{bclogo}[arrondi =0.1,   noborder = true, logo=\bcvelo, marge=4]{~\textbf{Activité} \textbf{\theactivite} {\itshape #1} }  \par}
{
\end{bclogo}
 \par \bigskip }


\newcounter{rque}
\newenvironment{remarque}
{\par \medskip    \addtocounter{rque}{1} \noindent  
\begin{bclogo}[arrondi =0.1,  ombre = true, barre=snake, noborder = true, logo=\bcinfo, marge=0]{~\textbf{Remarque} \textbf{\therque}}  \par }
{
\end{bclogo}
 \par \bigskip }

\newcounter{def}
\newenvironment{definition}[1]
{\par \medskip   \addtocounter{def}{1} \noindent  
\begin{bclogo}[arrondi =0.1,  ombre = true, barre=none, logo=\bcbook, marge=4]{~\textbf{Définition} \textbf{\thedef} {\itshape #1} }  \par}
{
\end{bclogo}
 \par \bigskip }
 
 
 \newcounter{cours}
\newenvironment{cours}[1]
{\par \medskip   \addtocounter{cours}{1} \noindent  
\begin{bclogo}[arrondi =0.1,  ombre = true, barre=none, logo=\bcbook, marge=4]{~\textbf{Point de cours} \textbf{\thecours} {\itshape #1} }  \par}
{
\end{bclogo}
 \par \bigskip }
 
 

\newenvironment{introduction}
{\par \medskip    \noindent  
 \begin {bclogo}[couleur = blue!5 , arrondi =0.1,logo=\bcrosevents, marge=4] {~\textbf{Introduction}    }
 \par }
{
\end{bclogo}
 \par \bigskip }
 
\newenvironment{memo}[1]
{\par \medskip    \noindent  
\begin{bclogo}[arrondi =0.1,  ombre = true, barre=none, logo=\bccle, marge=4]{~\textbf{À retenir}  {\itshape #1} }  \par}
{
\end{bclogo}
 \par \bigskip }
 
\newcounter{exple}
\newenvironment{exemple}[1]
{\par \medskip   \addtocounter{exple}{1} \noindent  
\begin{bclogo}[arrondi =0.1,   noborder = true, logo=\bclampe, marge=4]{~\textbf{Exemple} \textbf{\theexple} {\itshape #1} }  \par}
{
\end{bclogo}
 \par \bigskip }




\newcounter{alg}
\newenvironment{algorithme}[1]
{\par \medskip   \addtocounter{alg}{1} \noindent  
 \begin {bclogo}[noborder = true, barre=zigzag,logo=\bcpython, marge=4] {~\textbf{Algorithmique} \textbf{\thealg} {\itshape #1} }  \par}
{
\end{bclogo}
 \par \bigskip }

\newcounter{prog}
\newenvironment{programme}[1]
{\par \medskip   \addtocounter{prog}{1} \noindent  
 \begin {bclogo}[noborder = true, barre=zigzag,logo=\bcpython, marge=4] {~\textbf{Programme} \textbf{\theprog} {\itshape #1} }  \par  \bigskip}
{
\end{bclogo}
 \par \bigskip }
 
\newcounter{logi}
\newenvironment{logique}[1]
{\par \medskip   \addtocounter{logi}{1} \noindent  
 \begin {bclogo}[noborder = true, barre=zigzag,logo=\bclampe, marge=4] {~\textbf{Logique} \textbf{\thelogi} {\itshape #1} }  \par}
{
\end{bclogo}
 \par \bigskip }


\newenvironment{methode}[1]
{\par \medskip    \noindent  
 \begin {bclogo}[arrondi =0.1,logo=\bcoutil, marge=4,noborder = true] {~\textbf{Méthode}   {\itshape #1} }  \par}
{
\end{bclogo}
 \par \bigskip }


\newcounter{histo}
\newenvironment{histoire}[1]
{\par \medskip   \addtocounter{histo}{1} \noindent  
 \begin {bclogo}[couleur = blue!10 , arrondi =0.1,logo=\bchorloge, marge=4] {~\textbf{Histoire} \textbf{\thehisto} {\itshape #1} }  \par}
{
\end{bclogo}
 \par \bigskip }




%Environnement contenu pour un document présentant une progression annuelle
\newenvironment{contenu}
{\par \medskip   \begin {bclogo}[ noborder = true,logo=\bccrayon] \noindent {\large \textbf{Contenu de la séance}} \vspace*{10pt} \par  }
{\end{bclogo}  \par \medskip }

%Environnement programme pour un document présentant une progression annuelle
%\newenvironment{programme}
%{\par \medskip   \begin {bclogo}[ noborder = true, barre=zigzag,logo=\bcinfo] \noindent {\large \textbf{Programme officiel}} \vspace*{10pt} \par  }
%{\end{bclogo}  \par \medskip }

%Environnement programme pour un document présentant une progression annuelle
\newenvironment{ressource}
{\par \medskip   \begin {bclogo}[ noborder = true,logo=\bcbook] \noindent {\large \textbf{Ressources}}\\vspace*{10pt} \par }
{\end{bclogo}  \par \medskip }




%%%%%%%%%%%%%%%%%%Maths divers%%%%%%%%%%%%%%%%%%%%%%%%%
%%%%%%%%%%%%%Nombres%%%%%%%%%%%%%%%%

%Ensemble prive de...
%\newcommand{\prive}{\boi}%{\backslash}

%Ensembles de nombres%%%%%%%%%%%%%%%%%
\newcommand{\R}{\mathbb{R}}
\newcommand{\N}{\mathbb{N}}
\newcommand{\D}{\mathbb{D}}
\newcommand{\Z}{\mathbb{Z}}
\newcommand{\Q}{\mathbb{Q}}
%\newcommand{\C}{\mathbb{C}}
\newcommand{\df}{~\ensuremath{]0;+\infty[}~}
\newcommand{\K}{\mathbb{K}}

%%%%%%%%Arithmetique%%%%%%%%%%
%PGCD, PPCM
\newcommand{\PGCD}{\mathop{\rm PGCD}\nolimits}
\newcommand{\PPCM}{\mathop{\rm PPCM}\nolimits}

%Intervalles
\newcommand{\interoo}[2]{]#1\, ;\, #2[}
\newcommand{\Interoo}[2]{\left]#1\, ;\, #2\right[}
\newcommand{\interof}[2]{]#1\, ;\, #2]}
\newcommand{\Interof}[2]{\left]#1\, ;\, #2\right]}
\newcommand{\interfo}[2]{[#1\, ;\, #2[}
\newcommand{\Interfo}[2]{\left[#1\, ;\, #2\right[}
\newcommand{\interff}[2]{[#1\, ;\, #2]}
\newcommand{\Interff}[2]{\left[#1\, ;\, #2\right]}
%\newcommand\interentiers #1#2{[\! [#1\, ;\, #2]\! ]}
\newcommand{\interentiers}[2]{\llbracket #1\, ;\, #2\rrbracket}
%


%%%%%%%%%%%%%%Nombres complexes%%%%%

\newcommand{\ic}{\text{i}}
%\newcommand{\I}{\text{i}}
\newcommand{\im}[1]{\text{Im}\left(#1\right)}
\newcommand{\re}[1]{\text{Re}\left(#1\right)}
\newcommand{\Arg}[1]{\text{arg}\left(#1\right)}
\newcommand{\Mod}[1]{\left[#1\right]}
%Parties entière, réelle, imaginaire, nombre i
\newcommand{\ent}[1]{\text{E}\left(#1\right)}
\renewcommand{\Re}{\mathop{\rm Re}\nolimits}
\renewcommand{\Im}{\mathop{\rm Im}\nolimits}
\renewcommand{\i}{\textrm{i}}

%%%%%%%%%%%Probabilites et statistiques%%%%%
\newcommand{\loibinom}[2]{\mathcal{B}\left(#1\ ; \ #2 \right)}
\newcommand{\loinorm}[2]{\mathcal{N}\left(#1\ ; \ #2 \right)}
\newcommand{\loiexp}[1]{\mathcal{E}\left(#1\right)}
\newcommand{\proba}[1]{\mathbb{P}\big(#1\big)}
\newcommand{\probacond}[2]{\mathbb{P}_{#2}\big(#1\big)}
\newcommand{\esperance}[1]{\mathbb{E}\left(#1\right)}
\newcommand{\variance}[1]{\mathbb{V}\left(#1\right)}
\newcommand{\ecart}[1]{\sigma\left(#1\right)}
\newcommand{\dnormx}{\frac{1}{\sqrt{2\pi}} \text{e}^{-\frac{x^2}{2}}}
\newcommand{\dnormt}{\frac{1}{\sqrt{2\pi}} \text{e}^{-\frac{t^2}{2}}}

%Covariance
\newcommand{\cov}{\mathop{\rm cov}\nolimits}
%


%%%%%%%%%%Analyse%%%%%%%%%%%

%%%%%%%%%%%Courbe%%%%%%%%%%%%
\newcommand{\courbe}[1]{\ensuremath{\mathcal{C}_{#1}}}

%%%%%%%Fonction exponentielle%%%%%
\newcommand{\fe}{~fonction exponentielle~}
\newcommand{\e}{\text{e}}

%Fonction cotangente
\newcommand{\cotan}{\mathop{\rm cotan}\nolimits}
%%%%%%%%%%%%%%%%%%%%%%%%%%%%%%%%%%%%%%%%%
%
%Fonctions hyperboliques
\newcommand{\ch}{\mathop{\rm ch}\nolimits}
\newcommand{\sh}{\mathop{\rm sh}\nolimits}


%%%%%%%%%%%%%%Limites%%%%%%
\newcommand{\limite}[2]{\lim\limits
_{x \to #1} #2}
\newcommand{\limitesuite}[1]{\lim\limits
_{n \to +\infty} #1}
\newcommand{\limiteg}[2]{\lim\limits
_{\substack{x \to #1 \\ x < #1 }} #2}
\newcommand{\limited}[2]{\lim\limits
_{\substack{x \to #1 \\ x > #1 }} #2}

%%%%%%%%%%Continuité%%%%%%%%%%%
\newcommand{\TVI}{théorème des valeurs intermédiaires}

%%%%%%%%%%%Suites%%%%%%%%%%%%
\newcommand{\suite}[1]{\ensuremath{\left(#1_{n}\right)}}
\newcommand{\Suite}[2]{\ensuremath{\left(#1\right)_{#2}}}
%

%%%%%%%%%%%%%%%Calcul intégral%%%%%%
\newcommand{\dx}{\ensuremath{\text{d}x}}		% dx
\newcommand{\dt}{\ensuremath{\text{d}t}}		% dt
\newcommand{\dtheta}{\ensuremath{\text{d}\theta}}		% dtheta
\newcommand{\dy}{\ensuremath{\text{d}y}}		% dy
\newcommand{\dq}{\ensuremath{\text{d}q}}		% dq

%%%Intégrale%%%
\newcommand{\integralex}[3]{\int_{#1}^{#2} #3 \ \dx}
\newcommand{\integralet}[3]{\int_{#1}^{#2} #3 \ \dt}
\newcommand{\integraletheta}[3]{\int_{#1}^{#2} #3 \ \dtheta}

%%%%%Equivalent%%
\newcommand{\equivalent}[1]{\build\sim_{#1}^{}}

%o et O%%%%
\renewcommand{\o}[2]{\build o_{#1\to #2}^{}}
\renewcommand{\O}[2]{\build O_{#1\to #2}^{}}



%%%%%%%%%%%%%%%Geometrie%%%%%%%%%%%%%%%%%%%%%%%

%%%%%%%%%%%%%%%Reperes%%%%%%%%%%%%%%
\def\Oij{\ensuremath{\left(\text{O},~\vect{\imath},~\vect{\jmath}\right)}}
\def\Oijk{\ensuremath{\left(\text{O},~\vect{\imath},~ \vect{\jmath},~ \vect{k}\right)}}
\def\Ouv{\ensuremath{\left(\text{O},~\vect{u},~\vect{v}\right)}}
\renewcommand{\ij}{(\vec\imath\, ;\vec\jmath\,)}
\newcommand{\ijk}{(\vec\imath\, ;\vec\jmath\, ;\vec k\,)}
\newcommand{\OIJ}{(O\,;\, I\,;\, J\,)}
\newcommand{\repere}[3]{\big(#1\, ;\,\vect{#2} ;\vect{#3}\big)}
\newcommand{\reperesp}[4]{\big(#1\, ;\,\vect{#2} ;\vect{#3} ;\vect{#4}\big)}

%%%%%%%%%Coordonnees%%%%%%%%%%%%%%
\newcommand{\coord}[2]{(#1\, ;\, #2)}
\newcommand{\bigcoord}[2]{\big(#1\, ;\, #2\big)}
\newcommand{\Coord}[2]{\left(#1\, ;\, #2\right)}
\newcommand{\coordesp}[3]{(#1\, ;\, #2\, ;\, #3)}
\newcommand{\bigcoordesp}[3]{\big(#1\, ;\, #2\, ;\, #3\big)}
\newcommand{\Coordesp}[3]{\left(#1\, ;\, #2\, ;\, #3\right)}
\newcommand{\Vcoord}[3]{\begin{pmatrix} #1 \\ #2 \\ #3 \end{pmatrix}}
%Symboles entre droites
%\newcommand{\paral}{\sslash}
\newcommand{\paral}{\mathop{/\!\! /}}
%

%%%%%%%%%Produit scalaire, Angles%%%%%%%%%%
\newcommand{\scal}[2]{\vect{#1} \, \cdot \, \vect{#2}}
\newcommand{\Angle}[2]{\left(\vect{#1} \, , \, \vect{#2}\right)}
\newcommand{\Anglegeo}[2]{\left(\widehat{\vect{#1} \, ; \, \vect{#2}}\right)}
\renewcommand{\angle}[1]{\widehat{#1}}
\newcommand{\anglevec}[2]{\left(\vect {#1}\, ,\,\vect {#2} \right)}
\newcommand{\anglevecteur}[2]{(#1\, , \, #2)}
\newcommand{\Anglevec}[2]{(\vecteur{#1}\, ,\,\vecteur{#2})}
\newcommand{\prodscal}[2]{#1 \, \cdot \, #2}
%


%Arc
%\newcommand{\arc}[1]{\wideparen{#1}}
\newcommand{\arcoriente}[1]{\overset{\curvearrowright}{#1}}
%
%


%%%%%%%%%%%%%%%Normes%%%%%%%%%%%%%%%%
\newcommand{\norme}[1]{\left\| #1\right\|}
\newcommand{\normebis}[1]{\delim{2pt}{\|}{9pt}\! #1\delim{2pt}{\|}{9pt}}
\newcommand{\normetriple}[1]{\left |\kern -.07em\left\| #1\right |\kern -.07em\right\|}
\newcommand{\valabs}[1]{\big| \, #1 \, \big|}
%

%%%%%%%%%%%%%%%%%%%%%%%%%%%Degré%%%%%%
%\newcommand{\Degre}{\ensuremath{^\circ}}
%La commande \degre est déjà définie dans le package babel

%%%%%%%%%%Vecteurs%%%%%%%%%%%
\newcommand{\vect}[1]{\mathchoice%
{\overrightarrow{\displaystyle\mathstrut#1\,\,}}%
{\overrightarrow{\textstyle\mathstrut#1\,\,}}%
{\overrightarrow{\scriptstyle\mathstrut#1\,\,}}%
{\overrightarrow{\scriptscriptstyle\mathstrut#1\,\,}}}



%%%%%%%%%%%%%Algebre%%%%%%%%%%%%%%%


%%%%%%%%%%Systemes%%%%%%%%%%%
%Systemes
\newcommand{\sys}[2]{
\left\lbrace
 \begin{array}{l}
  \negthickspace\negthickspace #1\\
  \negthickspace\negthickspace #2\\
 \end{array}
\right.\negthickspace\negthickspace}
\newcommand{\Sys}[3]{
\left\lbrace
 \begin{array}{l}
  #1\\
  #2\\
  #3\\
 \end{array}
\right.}
\newcommand{\Sysq}[4]{
\left\lbrace
 \begin{array}{l}
  #1\\
  #2\\
  #3\\
  #4\\
 \end{array}
\right.}
%
%

%%%%%%%%%%%%%%%%Matrices%%%%%%%%%%%%%%%%%%
%Comatrice
\newcommand{\com}{\mathop{\rm com}\nolimits}
%
%
%Trace
\newcommand{\tr}{\mathop{\rm tr}\nolimits}
%
%
%Transposee
\newcommand{\transposee}[1]{{\vphantom{#1}}^t\negmedspace #1}
%
%
%Noyau
\newcommand{\Ker}{\mathop{\rm Ker}\nolimits}
%
%

%
%Matrices
\newcommand{\Mn}{\mathcal M_n}
\newcommand{\matrice}[4]{
\left(
 \begin{array}{cc}
  #1 & #2 \\
  #3 & #4
 \end{array}
\right)}

\newcommand{\Matrice}[9]{
\left(
 \begin{array}{ccc}
  #1 & #2 & #3\\
  #4 & #5 & #6\\
  #7 & #8 & #9
 \end{array}
\right)}
\newcommand{\Vect}[3]{
\left(\negmedspace
 \begin{array}{c}
  #1\\
  #2\\
  #3
 \end{array}\negmedspace
\right)}
\newcommand{\Ideux}{\matrice{1}{0}{0}{1}}
\newcommand{\Itrois}{\Matrice{1}{0}{0}{0}{1}{0}{0}{0}{1}}
%
%
%Determinants
\newcommand{\determinant}[4]{
\left|
 \begin{array}{cc}
  #1 & #2 \\
  #3 & #4
 \end{array}
\right|}
\newcommand{\Determinant}[9]{
\left|
 \begin{array}{ccc}
  #1 & #2 & #3\\
  #4 & #5 & #6\\
  #7 & #8 & #9
 \end{array}
\right|}

\begin{document}
\maketitle

\renewcommand*\contentsname{Table des matières}
{
\hypersetup{linkcolor=}
\setcounter{tocdepth}{3}
\tableofcontents
}
\hypertarget{cruxe9dits}{%
\section*{Crédits}\label{cruxe9dits}}
\addcontentsline{toc}{section}{Crédits}

\emph{Ce cours est inspiré du chapitre 24 du manuel NSI de la collection
Tortue chez Ellipse, auteurs : Ballabonski, Conchon, Filliatre, N'Guyen.
J'ai également consulté
\href{https://gitlab.com/lyceeND/1ere/-/blob/master/2019_20/2_Systeme_1/systeme19.md}{le
cours de Guillaume Connan}, le livre
\href{https://laurentbloch.net/MySpip3/IMG/pdf/histsys-screen-20200727.pdf}{Splendeurs
et servitudes des Systèmes d'exploitation Histoire, fonctionnement,
enjeux} de Laurent Bloch, le livre
\href{https://www.eyrolles.com/Informatique/Livre/la-ligne-de-commande-par-l-exemple-9782351410721/}{La
ligne de commande par l'exemple} de Vincent Fourmond et le livre
\href{https://www.eyrolles.com/Informatique/Livre/parlez-vous-shell--9782729877590/}{Parlez-vous
Shell ?} de Thomas Hugel.}

\hypertarget{chronologie-des-systuxe8mes-dexploitation}{%
\section{Chronologie des systèmes
d'exploitation}\label{chronologie-des-systuxe8mes-dexploitation}}

\begin{itemize}
\item
  1950 premiers calculateurs (ENIAC, EDSAC) branchements manuels pas de
  système d'exploitation
\item
  1956 GM-NAA I/O premier système d'exploitation sur IBM 704
\item
  1967 MULTICS développé aux Bell Labs premier système d'exploitation à
  temps partagé
\item
  1970 UNIX développé aux Bell Labs premier système d'exploitatio
  multitâches et multiutilisateurs, écrit en langage C
\item
  1978 BSD variante d'UNIX dévéloppée à l'université de Berkeley,
  ancêtre de MacOS
\item
  1980 les ordinateurs personnels d'IBM utilisent le système
  d'exploitation MS-DOS de Microsoft
\item
  1990 Microsoft développé Windows un système d'exploitation avec
  interface graphique inspiré du système disponible sur les machines
  Apple
\item
  1991 système d'exploitation libre Linux, noyau développé par Linus
  Torvalds, outils logiciels du projet GNU, distribué sous licence libre
\item
  2001 système d'exploitation MacOS sur les machines Apple, dérivé de
  BSD
\item
  2007 Apple développe un version de MacOS pour téléphones portables
\item
  2008 Google développe le système d'exploitation Android pour
  téléphones portables, basé sur le noyau libre Linux avec ajout de
  logiciels propriétaires
\end{itemize}

\hypertarget{duxe9finition-dun-systuxe8me-dexploitation}{%
\section{Définition d'un système
d'exploitation}\label{duxe9finition-dun-systuxe8me-dexploitation}}

Un \emph{système d'exploitation} est un logiciel, ou ensemble de
programmes, qui sert d'interface entre les programmes exécutés par
l'utilisateur et les ressources matérielles d'un ordinateur.

\begin{figure}
\centering
\includegraphics{https://upload.wikimedia.org/wikipedia/commons/thumb/e/ed/Operating_system_placement-fr.svg/250px-Operating_system_placement-fr.svg.png}
\caption{Illustration du rôle d'un système d'exploitation}
\end{figure}

Il est à la fois :

\begin{itemize}
\item
  une \emph{machine virtuelle} qui présente une interface simplifiée
  d'accès aux ressources (processeur, mémoire, périphériques
  d'entrée/sortie, réseau \ldots) pour les autres programmes et pour
  l'utilisateur
\item
  un \emph{chef d'orchestre} et un \emph{administrateur} :

  \begin{itemize}
  \tightlist
  \item
    c'est le premier programme exécuté au démarrage de l'ordinateur
  \item
    il gère l'accès concurrent aux ressources par les différents
    programmes (ordonnancement de l'utilisation du processeur par les
    programmes en cours d'exécution ou processus, sécurisation de la
    mémoire) ou utilisateurs (droits d'accès du système de fich
  \end{itemize}
\end{itemize}

Sitographie :

\begin{itemize}
\tightlist
\item
  \href{https://fr.wikipedia.org/wiki/Syst\%C3\%A8me_d\%27exploitation}{Page
  Wikipedia}
\item
  \href{https://www.lumni.fr/video/comprendre-ce-qu-est-un-systeme-d-exploitation}{Présentation
  en video}
\end{itemize}

\hypertarget{interface-utilisateur-dun-systuxe8me-dexploitation-le-shell}{%
\section{Interface utilisateur d'un système d'exploitation : le
shell}\label{interface-utilisateur-dun-systuxe8me-dexploitation-le-shell}}

Une interface entre l'utilisateur et le système d'exploitation s'appelle
un \href{https://fr.wikipedia.org/wiki/Shell_Unix}{shell} ou
\emph{interpréteur de commandes}.

Le rôle d'un \emph{shell} est de prendre une entrée de l'utilisateur, de
la traduire en instructions compréhensibles par le système
d'exploitation et de renvoyer la réponse du système à l'utilisateur.

Il existe deux grandes catégories de \emph{shell} :

\begin{itemize}
\tightlist
\item
  les \emph{interfaces textuelles} comme
  \href{https://fr.wikipedia.org/wiki/Bourne-Again_shell}{bash}, le plus
  commun sur les systèmes de la famille
  \href{https://fr.wikipedia.org/wiki/Bourne-Again_shell}{UNIX}.
\item
  les \emph{interfaces graphiques} qu'on retrouve dans les systèmes
  d'exploitation grand public tels que
  \href{https://fr.wikipedia.org/wiki/Microsoft_Windows}{Windows}
\end{itemize}

Nous utiliserons
\href{https://fr.wikipedia.org/wiki/Bourne-Again_shell}{bash} dans les
TP. La syntaxe d'une commande
\href{https://fr.wikipedia.org/wiki/Bourne-Again_shell}{bash} est simple
: le nom de la commande peut être suivi d'options facultatives
introduites par un tiret et d'arguments :
\texttt{nom\_commande\ -option1\ -option2\ argument1\ argument2}.

Par exemple, la commande ci-dessous demande d'utiliser la commande
\texttt{ls} pour afficher les informations détaillées (option
\texttt{-l}) sur le fichier \texttt{exemple}:

\begin{verbatim}
junier@fredportable:~/bac$ ls -l exemple
-rw-rw-r-- 1 nobody mail 0 août  16 12:05 exemple
\end{verbatim}

Chaque commande est un programme enregistré dans un fichier, par exemple
\texttt{/usr/bin/ls} pour la commande \texttt{ls}:

\begin{verbatim}
junier@fredportable:~/bac$ which ls
/usr/bin/ls
junier@fredportable:~/bac$ ls -l /usr/bin/ls
-rwxr-xr-x 1 root root 142144 sept.  5  2019 /usr/bin/ls
\end{verbatim}

Enfin chaque commande possède trois flux d'entrée / sortie :

\begin{itemize}
\tightlist
\item
  \emph{l'entrée standard} qui est par défaut le clavier (les options et
  arguments sont traités puis transmis par
  \href{https://fr.wikipedia.org/wiki/Bourne-Again_shell}{bash} sur
  l'entrée standard de la commande)
\item
  \emph{la sortie standard} qui est par défaut l'écran
\item
  \emph{la sortie d'erreur} qui est par défaut l'écran également
\end{itemize}

Il existe des opérateurs \texttt{\textgreater{}} et \texttt{\textless{}}
pour rediriger les flux d'entrée / sortie.

Considérons un exemple . La commande \texttt{ls} sans argument affiche
le contenu du répertoire courant sur la sortie standard et la commande
\texttt{cat} concatène les fichiers fournis en paramètres ou sur son
entrée standard et renvoie le résultat sur sa sortie standard :

\begin{verbatim}
    junier@fredportable:~/bac$ ls 
    entree
    junier@fredportable:~/bac$ cat --help 
    utilisation : cat [OPTION] [FICHIER]...
    Concaténer les FICHIERs vers la sortie standard.
    ............
    junier@fredportable:~/bac$ cat entree
    contenu de entree
    junier@fredportable:~/bac$ cat entree > sortie
    junier@fredportable:~/bac$ ls
    entree  sortie
    junier@fredportable:~/bac$ cat sortie
    contenu de entree
\end{verbatim}

Nous fournirons un memento des principales commandes de
\href{https://fr.wikipedia.org/wiki/Bourne-Again_shell}{bash} communes à
la plupart des \emph{shell} des systèmes de la famille
\href{https://fr.wikipedia.org/wiki/Bourne-Again_shell}{UNIX}.

Un memento en ligne : \url{https://juliend.github.io/linux-cheatsheet/}

\hypertarget{attributs-du-systuxe8me-dexploitation-le-mode-noyau}{%
\section{Attributs du système d'exploitation : le mode
noyau}\label{attributs-du-systuxe8me-dexploitation-le-mode-noyau}}

Les programmes ne peuvent pas accéder directement aux ressources
matérielles (mémoire, processeur, périphériques d'entrée/sortie) sinon
il y aurait des conflits : par exemple deux programmes pourraient écrire
dans la même zone mémoire.

Le système d'exploitation possède un mode d'exécution privilégiée : le
\emph{mode noyau}, qui lui donne un accès unique et total aux
ressources. Les autres programmes s'exécutent en \emph{mode normal}.
Pour accéder aux ressources, ils font appel au système d'exploitation à
l'aide de primitives qu'on nomme \emph{appels systèmes}. Pour les
systèmes d'exploitation dérivés d'\href{}{UNIX}, ces appels systèmes
sont normalisés par le standard
\href{https://fr.wikipedia.org/wiki/POSIX}{POSIX} pour Portable
Operating System Interface.

Le \textbf{noyau} désigne l'ensemble des programmes au coeur du système
d'exploitation, qui s'exécutent en particulier au démarrage de
l'ordinateur. Par exemple, le noyau du système d'exploitation libre
Linux est téléchargeable sur \url{https://www.kernel.org/}.

\hypertarget{attributs-du-systuxe8me-dexploitation-le-contruxf4le-de-lactivituxe9-des-programmes}{%
\section{Attributs du système d'exploitation : le contrôle de l'activité
des
programmes}\label{attributs-du-systuxe8me-dexploitation-le-contruxf4le-de-lactivituxe9-des-programmes}}

Un programme ne peut être exécuté que si le système d'exploitation
l'autorise. L'exécution d'un programme peut être interrompue par le
système d'exploitation, pour libérer une ressource ou en cas d'erreur.
On parle d'\emph{interruption système}.

Un programme en cours d'exécution s'appelle un \emph{processus} et le
système d'exploitation arbitre le partage de temps de calcul du
processeur par les processus concurrents. Il joue un rôle
\emph{d'ordonnanceur} pour éviter par exemple qu'un processus s'accapare
le processeur.

Le système d'exploitation maintient une base de temps unique pour tous
les processus et fournit des services de gestion du temps :
estampillage, chronologie, attente, réveil.

La commande de \emph{shell} \texttt{ps} permet d'afficher la liste des
processus : l'option \texttt{U} permet de sélectionner l'utilisateur qui
a lancé les processus. S'il s'agit d'un processus qu'on a lancé ou si on
est superutilisateur, peut alors ``tuer'' le processus avec la commande
\texttt{kill} suivie du numéro du processus.

\begin{verbatim}
    junier@fredportable:~$ ps U
    PID TTY      STAT   TIME COMMAND
    13845 pts/5    S+     0:00 python3
    13900 pts/4    R+     0:00 ps U junier
    junier@fredportable:~$ kill 13845
\end{verbatim}

\hypertarget{attributs-du-systuxe8me-dexploitation-le-contruxf4le-daccuxe8s-aux-ressources}{%
\section{Attributs du système d'exploitation : le contrôle d'accès aux
ressources}\label{attributs-du-systuxe8me-dexploitation-le-contruxf4le-daccuxe8s-aux-ressources}}

Le système d'exploitation est le seul à pouvoir contrôler l'accès aux
ressources : processeur, mémoire, périphériques d'entrée/sortie.

La gestion et la sécurisation de la mémoire est une tâche importante,
avec un système de \emph{mémoire virtuelle} qui donne à chaque programme
l'illusion d'accéder à toute la mémoire.

Des algorithmes permettent d'éviter les situations
d'\href{https://fr.wikipedia.org/wiki/Interblocage}{interblocage ou
étreinte fatale}.

Un exemple concret d'interblocage peut se produire lorsque deux
processus essayent d'acquérir deux ressources dans un ordre différent.
Par exemple avec deux ressources (M1 et M2), deux processus (P1 et P2)
et la séquence d'opération suivante :

\begin{verbatim}
P1 acquiert M1.
P2 acquiert M2.
P1 attend pour acquérir M2 (qui est détenu par P2).
P2 attend pour acquérir M1 (qui est détenu par P1).
\end{verbatim}

Dans cette situation, les deux processus sont définitivement bloqués.

La commande \emph{shell} \texttt{free} permet d'afficher l'utilisation
de la mémoire, l'option \texttt{-m} convertit les mesures en megaoctets
(Mo). Dans l'exemple ci-dessous, la mémoire vive disponible est de 3856
Mo, 618 Mo sont utilisés pour les buffers/cache et 254 sont libres soit
un total de 3856 - 618 - 254 = 2984 Mo utilisés. La seconde ligne
concerne un espace du disque dur, nommé
\href{https://doc.ubuntu-fr.org/swap}{swap}, utilisé par le système pour
étendre la mémoire vive en mémoire virtuelle.

\begin{verbatim}
junier@fredportable:~/bac$ free -m
            total       utilisé      libre     partagé tamp/cache   disponible
Mem:           3856        2983         254         302         618         334
Partition d'échange:        1303        1288          14
\end{verbatim}

\hypertarget{attributs-du-systuxe8me-dexploitation-le-contruxf4le-du-systuxe8me-de-fichiers}{%
\section{Attributs du système d'exploitation : le contrôle du système de
fichiers}\label{attributs-du-systuxe8me-dexploitation-le-contruxf4le-du-systuxe8me-de-fichiers}}

\hypertarget{systuxe8me-de-fichiers}{%
\subsection{Système de fichiers}\label{systuxe8me-de-fichiers}}

Sur les supports de mémoire persistants (disques durs, clefs USB\ldots),
les informations sont regroupées par le système d'exploitation dans des
\emph{fichiers} qui sont organisés à travers un
\href{https://fr.wikipedia.org/wiki/Syst\%C3\%A8me_de_fichiers}{système
de fichiers}.

Dans les systèmes d'exploitation qui respectent le standard
\href{https://fr.wikipedia.org/wiki/POSIX}{POSIX}, on distingue
plusieurs catégories de fichiers :

\begin{itemize}
\tightlist
\item
  \emph{fichiers réguliers} ou \emph{fichiers textes} qui contiennent
  des suites de caractères et qui sont lisibles par des humains
\item
  \emph{fichiers binaires} qui sont des suites d'octets non lisibles par
  des humains
\item
  \emph{fichiers exécutables} qui sont des programmes, ils peuvent être
  des fichiers textes ou binaires.
\item
  \emph{répertoires} qui sont des listes de fichiers : ils servent de
  conteneur à fichiers.
\end{itemize}

Les \emph{répertoires} pouvant contenir d'autres fichiers, un système de
fichiers \href{https://fr.wikipedia.org/wiki/POSIX}{POSIX} possède une
structure hiérarchique qui est une \emph{arborescence} avec un
\emph{répertoire racine} symbolisé par un /.

\begin{figure}
\centering
\includegraphics{https://www.frederic-junier.org/ISN/Architecture-HTML-CSS-Internet-2019/pages/images/arborescence.png}
\caption{Un exemple d'arborescence}
\end{figure}

Les fichiers sont repérés par leur \emph{chemin} :

\begin{itemize}
\item
  \emph{chemin absolu} depuis la racine : par exemple si un fichier
  \texttt{exemple.txt} se trouve dans le répertoire \texttt{sandbox},
  contenu dans le répertoire \texttt{maurice}, lui-même contenu dans le
  répertoire \texttt{home}, lui-même dans le répertoire racine, son
  \emph{chemin absolu} est \texttt{/home/maurice/sandbox/exemple.txt}
\item
  ou \emph{chemin relatif} qui est relatif au \emph{répertoire courant}
  ou \emph{répertoire de travail} où l'utilisateur se trouve. Par
  exemple, si le répertoire courant est \texttt{maurice}, le
  \emph{chemin relatif} du fichier précédent est
  \texttt{sandbox/exemple.txt}.
\end{itemize}

\hypertarget{gestion-des-droits-et-permissions}{%
\subsection{Gestion des droits et
permissions}\label{gestion-des-droits-et-permissions}}

Le système d'exploitation gère les \emph{droits et permissions} sur les
fichiers (dont les répertoires) pour les différents \emph{utilisateurs}.

\hypertarget{utilisateurs-et-permissions}{%
\subsubsection{Utilisateurs et
permissions}\label{utilisateurs-et-permissions}}

Pour un fichier, on distingue :

\begin{itemize}
\item
  trois types d'utilisateurs :

  \begin{itemize}
  \tightlist
  \item
    le \emph{propriétaire} (ou \emph{owner}) noté \texttt{u}
  \item
    le \emph{groupe principal} (ou \emph{group}) noté \texttt{g}
  \item
    un \emph{autre utilisateur} (ou \emph{other}) noté \texttt{o}
  \end{itemize}
\item
  trois types de permissions :

  \begin{itemize}
  \tightlist
  \item
    \emph{lecture} (caractère \emph{r} si attribué ou \emph{-} sinon)
  \item
    \emph{écriture} (caractère \emph{w} si attribué ou \emph{-} sinon)
  \item
    \emph{exécution} (caractère \emph{x} si attribué ou \emph{-} sinon)
  \end{itemize}
\end{itemize}

\begin{figure}
\centering
\includegraphics{https://upload.wikimedia.org/wikipedia/commons/0/0e/Unix-file-rights.png}
\caption{droits unix}
\end{figure}

\emph{Source : Wikimedia Commons, Jimbotyson / CC BY-SA
(https://creativecommons.org/licenses/by-sa/3.0)}

Pour les répertoires, la signification des permissions est précisée dans
ce tableau. Attention, si on a le droit d'écriture sur un répertoire on
peut supprimer tous les fichiers qu'il contient même ceux dont on n'est
pas propriétaire !

\begin{longtable}[]{@{}lll@{}}
\toprule
& Type de permission & pour un répertoire\tabularnewline
\midrule
\endhead
r & lecture & lister le contenu\tabularnewline
w & écriture & ajouter, supprimer, renommer des fichiers\tabularnewline
x & exécution & entrer dedans\tabularnewline
\bottomrule
\end{longtable}

En général le nom d'utilisateur figure dans le prompt de l'interpréteur
de commandes / shell, mais la commande \texttt{id} permet de l'afficher
avec la liste des groupes auxquels il appartient :

\begin{verbatim}
junier@fredportable:~$ id
uid=1000(junier) gid=1000(junier) groupes=1000(junier),4(adm),8(mail),24(cdrom),27(sudo),30(dip),46(plugdev),120(lpadmin),131(lxd),132(sambashare)
junier@fredportable:~$ sudo apt install pandoc
[sudo] Mot de passe de junier :
\end{verbatim}

Certaines commandes, comme l'installation de logiciels avec un
gestionnaire de paquets comme \texttt{apt} ne sont exécutables qu'avec
les droits du \emph{superutilisateur} \texttt{root}. Si on appartient au
groupe \texttt{sudo} (groupe des administrateurs), on peut élever ses
droits et agir comme \texttt{root} en préfixant la commande de
\texttt{sudo}. et en saisissant son mot de passe utilisateur. Certains
systèmes permettent de devenir directement \texttt{root} avec la
commande \texttt{su} en saisissant le mot de passe de \texttt{root}.
\emph{L'utilisateur \texttt{root} a tous les droits sur le système de
fichiers.}

\hypertarget{exemples-de-lectures-de-permissions}{%
\subsubsection{Exemples de lectures de
permissions}\label{exemples-de-lectures-de-permissions}}

On peut afficher les droits d'un fichier avec la commande
\texttt{ls\ -l}, en considérant les 10 premiers caractères :

\begin{itemize}
\item
  le premier précise le type de fichier : \texttt{d} pour un répertoire
  et tiret \texttt{-} pour un autre fichier
\item
  les 9 caractères suivants dénotent les attributs pour les trois types
  de permission avec de gauche à droite le \emph{propriétaire}, le
  \emph{groupe}, les \emph{autres}.
\item
  Exemple 1 :

\begin{verbatim}
  junier@fredportable:~/sandbox$ ls -l histoire-systeme.txt         
  -rw-rw-r-- 1 junier junier 1136 août  15 13:56 histoire-systeme.txt
\end{verbatim}

  Dans cet exemple le fichier texte \texttt{histoire-systeme.txt} a pour
  propriétaire \texttt{junier}, pour groupe \texttt{junier} et pour
  droits :
\end{itemize}

\begin{longtable}[]{@{}llll@{}}
\toprule
& Lecture & Écriture & Exécution\tabularnewline
\midrule
\endhead
Propriétaire & oui & oui & non\tabularnewline
Groupe & oui & oui & non\tabularnewline
Autres & oui & non & non\tabularnewline
\bottomrule
\end{longtable}

\begin{itemize}
\item
  Exemple 2 :

\begin{verbatim}
  junier@fredportable:~/sandbox$ ls -l /bin/python3.8
  -rwxr-xr-x 1 root root 5453504 juil. 16 16:00 /bin/python3.8
\end{verbatim}

  Dans cet exemple le fichier binaire \texttt{/bin/python3.8} a pour
  propriétaire \texttt{root}, pour groupe \texttt{root} et pour droits :
\end{itemize}

\begin{longtable}[]{@{}llll@{}}
\toprule
& Lecture & Écriture & Exécution\tabularnewline
\midrule
\endhead
Propriétaire & oui & oui & oui\tabularnewline
Groupe & oui & non & oui\tabularnewline
Autres & oui & non & oui\tabularnewline
\bottomrule
\end{longtable}

\begin{itemize}
\item
  Exemple 3 : pour afficher les informations détaillées d'un répertoires
  il faut rajouter l'option \texttt{-d} :

\begin{verbatim}
  junier@fredportable:/var$ ls -l -d mail
  drwxrwsr-x 2 root mail 4096 avril 23 09:32 mail
\end{verbatim}
\end{itemize}

Dans cet exemple le dossier \texttt{/var/mail} a pour propriétaire
\texttt{root}, pour groupe \texttt{mail} et pour droits :

\begin{longtable}[]{@{}llll@{}}
\toprule
& Lecture & Écriture & Exécution\tabularnewline
\midrule
\endhead
Propriétaire & oui & oui & oui\tabularnewline
Groupe & oui & oui & oui\tabularnewline
Autres & oui & non & oui\tabularnewline
\bottomrule
\end{longtable}

On peut noter que pour le groupe, le droit d'exécution est positionné à
\texttt{s}, qui correspond au
\href{https://fr.wikipedia.org/wiki/Permissions_UNIX\#Droit_SGID}{sgid}
: tout fichier créé appartient au groupe propriétaire du répertoire.

Attention, si tout membre du groupe \texttt{mail} peut écrire dans le
répertoire \texttt{var}, il peut aussi supprimer un fichier créé par un
autre utilisateur même s'il s'agit du superutilisateur \texttt{root} :

\begin{verbatim}
junier@fredportable:/var/mail$ ls -l
total 0
-rw-r--r-- 1 root mail 0 août  16 10:27 mail_important
junier@fredportable:/var/mail$ echo 'Bonjour' > mon_message
junier@fredportable:/var/mail$ ls -l
total 4
-rw-r--r-- 1 root   mail 0 août  16 10:27 mail_important
-rw-rw-r-- 1 junier mail 8 août  16 10:28 mon_message
junier@fredportable:/var/mail$ rm mail_important
rm : supprimer 'mail_important' qui est protégé en écriture et est du type « fichier vide » ? y
junier@fredportable:/var/mail$ ls -l
total 4
-rw-rw-r-- 1 junier mail 8 août  16 10:28 mon_message
\end{verbatim}

\hypertarget{modification-des-droits-permissions}{%
\subsubsection{Modification des droits /
permissions}\label{modification-des-droits-permissions}}

\begin{itemize}
\item
  Seul le propriétaire d'un fichier (ou répertoire) ou l'utilisateur
  \texttt{root} peuvent modifier les permissions d'un fichier (ou
  répertoire) avec la commande \texttt{chmod}.

  La syntaxe est la suivante, on a noté entre crochets les options
  facultatives :

\begin{verbatim}
  chmod [ugo][+-=][rwx] fichier
\end{verbatim}

  Voici un exemple commenté :

  \begin{itemize}
  \item
    Lorsqu'on crée un fichier, par défaut il a les droits de lecture et
    d'écriture pour l'utilisateur et son groupe principal et de lecture
    pour les autres.

\begin{verbatim}
  junier@fredportable:~/bac$ touch fichier
  junier@fredportable:~/bac$ ls -l
  total 0
  -rw-rw-r-- 1 junier junier 0 août  16 11:05 fichier
  junier@fredportable:~/bac$ echo "echo exécution" > fichier
  junier@fredportable:~/bac$ cat fichier 
  echo exécution
\end{verbatim}
  \item
    On ajoute le droite d'exécution à \texttt{fichier} pour son
    propriétaire et on vérifie que la commande qu'il contient est bien
    exécutée lorsqu'on écrit \texttt{./fichier} :

\begin{verbatim}
  junier@fredportable:~/bac$ chmod u+x fichier
  junier@fredportable:~/bac$ ./fichier 
  exécution
  junier@fredportable:~/bac$ ls -l
  total 4
  -rwxrw-r-- 1 junier junier 16 août  16 11:06 fichier
\end{verbatim}
  \item
    On enlève le droit d'écriture au propriétaire et on vérifie qu'on ne
    peut plus écrire dans `fichier

\begin{verbatim}
  junier@fredportable:~/bac$ chmod u-w fichier 
  junier@fredportable:~/bac$ echo "echo ajout" > fichier
  bash: fichier: Permission non accordée
  junier@fredportable:~/bac$ ls -l
  total 4
  -r-xrw-r-- 1 junier junier 16 août  16 11:06 fichier
\end{verbatim}
  \item
    On ajoute le droit d'écriture au propriétaire et aux autres :

\begin{verbatim}
  junier@fredportable:~/bac$ chmod uo+w fichier
  junier@fredportable:~/bac$ ls -l
  total 4
  -rwxrw-rw- 1 junier junier 16 août  16 11:06 fichier
\end{verbatim}
  \item
    On enlève le droit de lecture au propriétaire, à son groupe et aux
    autres :

\begin{verbatim}
  junier@fredportable:~/bac$ chmod ugo-r fichier
  junier@fredportable:~/bac$ ls -l
  total 4
  --wx-w--w- 1 junier junier 16 août  16 11:06 fichier
\end{verbatim}
  \item
    Pour remettre ce droit de lecture à tous, on peut utiliser le
    raccourci \texttt{a} au lieu de \texttt{ugo} :

\begin{verbatim}
  junier@fredportable:~/bac$ chmod a+r fichier 
  junier@fredportable:~/bac$ ls -l
  total 4
  -rwxrw-rw- 1 junier junier 16 août  16 11:06 fichier
\end{verbatim}
  \item
    On positionne tous les droits à \texttt{rwx} sur \texttt{fichier}
    pour tous les utilisateurs. On aurait pu écrire
    \texttt{chmod\ a=rwx\ fichier}.

\begin{verbatim}
  junier@fredportable:~/bac$ chmod ugo=rwx fichier 
  junier@fredportable:~/bac$ ls -l
  total 8
  -rwxrwxrwx 1 junier junier   16 août  16 11:06 fichier
\end{verbatim}
  \item
    On crée un répertoire nommé \texttt{repertoire} avec deux fichiers.
    On peut noter qu'un répertoire est toujours créé par défaut avec le
    droit d'exécution pour tous, car celui-ci permet (ou non) de
    traverser le répertoire.

\begin{verbatim}
  junier@fredportable:~/bac$ mkdir repertoire
  junier@fredportable:~/bac$ ls -l -d repertoire/
  drwxrwxr-x 2 junier junier 4096 août  16 11:10 repertoire/
  junier@fredportable:~/bac$ cd repertoire/
  junier@fredportable:~/bac/repertoire$ touch fichier2
  junier@fredportable:~/bac/repertoire$ touch fichier3
  junier@fredportable:~/bac/repertoire$ ls -l
  total 0
  -rw-rw-r-- 1 junier junier 0 août  16 11:10 fichier2
  -rw-rw-r-- 1 junier junier 0 août  16 11:12 fichier3
\end{verbatim}
  \item
    On modifie récursivement les droits sur le répertoire
    \texttt{repertoire} et tous ses fichiers avec l'option \texttt{-R}
    dans la commande \texttt{chmod\ -R\ ugo=rx\ repertoire}. On peut
    vérifier qu'on ne peut plus supprimer les fichiers contenus dans
    \texttt{repertoire}.

\begin{verbatim}
  junier@fredportable:~/bac$ chmod -R ugo=rx repertoire
  junier@fredportable:~/bac$ ls -l repertoire
  total 0
  -r-xr-xr-x 1 junier junier 0 août  16 11:10 fichier2
  -r-xr-xr-x 1 junier junier 0 août  16 11:12 fichier3
  junier@fredportable:~/bac$ ls -l -d repertoire/
  dr-xr-xr-x 2 junier junier 4096 août  16 11:12 repertoire/
  junier@fredportable:~/bac$ cd repertoire/
  junier@fredportable:~/bac/repertoire$ rm fichier2
  rm : supprimer 'fichier2' qui est protégé en écriture et est du type « fichier vide » ? y 
  rm: impossible de supprimer 'fichier2': Permission non accordée
\end{verbatim}
  \end{itemize}
\item
  Le superutilisateur \texttt{root} peut changer le propriétaire d'un
  fichier avec la commande \texttt{chown} ou son groupe principal avec
  la commande \texttt{chgrp} :

\begin{verbatim}
      junier@fredportable:~/bac$ ls -l fichier
      -rw-rw-r-- 1 nobody mail 0 août  16 12:05 fichier
      junier@fredportable:~/bac$ sudo chown nobody fichier
      junier@fredportable:~/bac$ ls -l
      -rw-rw-r-- 1 nobody junier    0 août  16 12:05 fichier
      junier@fredportable:~/bac$ sudo chgrp mail  fichier
      junier@fredportable:~/bac$ ls -l
      -rw-rw-r-- 1 nobody mail      0 août  16 12:05 fichier
\end{verbatim}
\end{itemize}

\hypertarget{licence-logicielle-et-systuxe8me-dexploitation}{%
\section{Licence logicielle et système
d'exploitation}\label{licence-logicielle-et-systuxe8me-dexploitation}}

\emph{Une licence de logiciel est un contrat qui octroie un droit
d'usage non exclusif} et par lequel le titulaire des droits d'auteur du
logiciel définit les conditions dans lesquelles ce programme peut être
utilisé, diffusé ou modifié par l'usager.

\hypertarget{contexte-historique}{%
\subsection{Contexte historique}\label{contexte-historique}}

De ses débuts jusqu'à l'aube des années 1980, le monde de l'informatique
fut dominé par fabricants de \texttt{hardware} tels que IBM. Dans les
années 50 et 60, les ordinateurs étaient de gros calculateurs et
l'informatique concernait principalement des spécialistes (ingénieurs,
chercheurs). Les logiciels ou \texttt{software} existaient mais étaient
indissociables de la machine sur laquelle ils étaient installés par le
fabricant. L'interopérabilité des programmes, la normalisation des
formats de données n'étaient pas dans l'air du temps. D'ailleurs, le
constructeur ne facturait pas le logiciel qu'il livrait avec leur
machine et fournissait les codes sources aux utilisateurs afin qu'ils
puissent adapter les ressources de l'ordinateur à leur usage ou
rectifier certains bugs.

Dans les années 70, des projets de développement logiciels d'importance
émergèrent comme le projet \textbf{UNIX} d'un système d'exploitation
multiutilisateurs (plusieurs comptes utilisateurs mais un seul
utilisateur à la fosi alors que le système MULTICS avait pour ambition
de supporter plusieurs utilisateurs simultanés). Il fut d'abord
développé par Ken Thompson et Denis Ritchie (créateur du lagage C) dans
les laboratoires Bell. Cette entreprise livra le code source à
l'université de Berkeley dont les étudiants améliorèrent le projet qui
fut distribué sous le nom de \textbf{Berkeley Software Distribution}
(BSD).

A cette époque, les échanges de code sources entre développeurs étaient
naturels et ils furent accélérés par le déploiement du réseau Arpanet
d'interconnexion des universités et centres de recherche, ancêtre de
l'internet.

L'essor de la micro-informatique dans les années 80, le changement de
profil des usagers (plus forcément des techniciens), firent du logiciel
le facteur dominant du business informatique. Les sociétés d'édition de
logiciels comme Microsoft supplantèrent les fabricants de hardware comme
IBM. Mais la protection de leurs savoir-faire et de leur code ressource
devint un facteur essentiel de leur stratégie de développement. L'époque
des pionniers de l'informatique qui partageaient leurs codes sources et
leurs idées était révolue.

En 1984, \textbf{Richard Stallman} démissionnait du MIT, crée la
\textbf{Free Software Foundation} (FSF) et lançait le projet
\textbf{GNU} de développement d'un système d'exploitation libre dont
aucun élément ne serait soumis à une \textbf{licence propriétaire} (la
première distribution Unix sans code propriétaire ne fut publiée qu'en
1989). Stallman créa la première \textbf{licence logicielle libre} : la
General Public License (GPL). Le projet GNU aboutit en 1991 au système
d'exploitation libre \textbf{GNU/Linux} après incorporation du dernier
élément manquant, le noyau (ou kernel), apporté par Linus Torvalds.

Des références :

\begin{itemize}
\item
  \url{http://www.gnu.org/gnu/gnu.html}
\item
  \url{http://fr.wikipedia.org/wiki/Unix}
\end{itemize}

\begin{figure}
\centering
\includegraphics{https://upload.wikimedia.org/wikipedia/commons/thumb/2/26/Carte_conceptuelle_du_logiciel_libre.svg/1024px-Carte_conceptuelle_du_logiciel_libre.svg.png}
\caption{Carte du libre}
\end{figure}

\emph{Source : Wikimedia Commons, René Mérou et autres auteurs}

\hypertarget{les-diffuxe9rents-types-de-licences-logicielles}{%
\subsection{Les différents types de licences
logicielles}\label{les-diffuxe9rents-types-de-licences-logicielles}}

On distingue :

\begin{itemize}
\item
  les \textbf{licences propriétaires} :\\
  Un éditeur concède à titre non exclusif un droit d'usage sur un
  logiciel dont il conserve les droits de propriété intellectuelle. Un
  \textbf{CLUF} (Contrat Licence Utilisateur Final) délimite les
  conditions d'usage du logiciel : limitation du nombre de copies, de
  postes sur lequel le logiciel peut être installé, interdiction de
  modification et de redistribution . En général, le code source n'est
  pas accessible (on parle de logiciel fermé) \ldots.
\item
  les \textbf{licences libres} ou \textbf{open-source} :\\
  elles respectent les quatre libertés fondamentales du logiciel libre :

  \begin{enumerate}
  \def\labelenumi{\arabic{enumi}.}
  \item
    \textbf{liberté d'exécuter le programme, pour tous les usages.}
  \item
    \textbf{liberté d'étudier le fonctionnement du programme, et de le
    modifier pour l'adapter à ses besoins}
  \item
    \textbf{liberté de redistribuer des copies.}
  \item
    \textbf{liberté de redistribuer aux autres des copies de versions
    modifiées}
  \end{enumerate}

  Les libertés 2 et 3 nécessitent l'accès au code source.\\
  Il ne faut pas confondre licence libre et domaine public : mettre un
  logiciel sous licence libre ne signifie pas abandonner ses droits
  d'auteurs (droit moral inaliénable) mais donner des permissions aux
  usagers qui vont au-delà de ce qu'autorise le droit d'auteur par
  défaut.
\end{itemize}

Attention à \emph{ne pas confondre logiciels libres et gratuits} (double
sens du mot free en Anglais):

\begin{itemize}
\item
  un \textbf{freeware} désigne usuellement un logiciel distribué
  gratuitement, indépendamment de sa licence d'utilisation. Le code
  source n'est pas forcément fourni ; dans ce cas ce n'est pas un
  logiciel libre.
\item
  un \textbf{shareware} désigne un logiciel distribué gratuitement et
  librement pendant une durée ou un nombre d'utilisations qui sont
  fixées par l'auteur. Au delà de ce délai il faut payer des royalties
  et le code source n'est pas fourni donc ce n'est pas un logiciel
  libre.
\end{itemize}

Références :

\begin{itemize}
\item
  \url{http://www.cndp.fr/savoirscdi/societe-de-linformation/cadre-reglementaire/le-coin-du-juriste/logiciel-libre-et-gratuiciel.html}
\item
  \url{http://www.gnu.org/philosophy/categories.html} et
  \url{http://www.gnu.org/philosophy/philosophy.html}
\end{itemize}

\hypertarget{licences-de-logiciels-libres}{%
\subsection{Licences de logiciels
libres}\label{licences-de-logiciels-libres}}

On distingue deux types de licences libres :

\begin{enumerate}
\def\labelenumi{\arabic{enumi}.}
\item
  Les \textbf{licences avec copyleft} :

  Le copyleft ( gauche d'auteur , jeu de mots sur copyright), n'est pas
  l'abandon du copyright (droits patrimoniaux) mais au contraire
  s'appuie dessus pour rendre libre un programme (respect des 4 libertés
  fondamentales) en obligeant toutes les versions modifiées à être
  libres également et redistribuées sous la même licence. Ainsi le
  logiciel (ou l'oeuvre numérique) sera éternellement libre. La licence
  \texttt{GPL} (General Public License) créée par Richard Stallman en
  1989 fut la première licence libre avec copyleft. La licence
  \texttt{CeCill} créée par l'INRIA est une adaptation de la licence GPL
  au droit français. La licence \texttt{Creative\ Commons} CC BY-SA 3.0,
  paternité et partage à l'identique, est une licence libre avec
  copyleft pour les oeuvres numériques (image, video \ldots).

  Il existe des formes plus ou moins fortes de copyleft. Ainsi si on
  inclut du code sous license GPL dans un autre, le programme hybride
  créé est contaminé et doit être distribué sous licence GPL. La license
  \textbf{LGPL} présente un copyleft plus faible : du code sous cette
  license peut être inclus dans un code sous une autre license sans que
  le programme hybride soit nécessairement sous license LGPL. C'est
  utile pour développer des bibliothèques de fonctions.
\item
  Les \textbf{licences sans copyleft} :

  L'utilisateur peut redistribuer et modifier le logiciel, mais aussi
  ajouter des restrictions et distribuer une version modifiée sous
  licence propriétaire. On parle aussi de licence permissive . La
  première licence libre sans copyleft fut la licence \texttt{BSD} pour
  le distributions de logiciels créés par l'université de Berkeley.
  \emph{La licence de Python est une licence libre sans copyleft}. Ce
  sont des licences compatibles avec la GPL mais la réciproque est
  fausse.
\end{enumerate}

\hypertarget{exemple-des-systuxe8mes-dexploitation}{%
\subsection{Exemple des systèmes
d'exploitation}\label{exemple-des-systuxe8mes-dexploitation}}

\begin{figure}
\centering
\includegraphics{https://upload.wikimedia.org/wikipedia/commons/thumb/7/77/Unix_history-simple.svg/1024px-Unix_history-simple.svg.png}
\caption{Arbre généalogique des systèmes UNIX}
\end{figure}

\begin{figure}
\centering
\includegraphics{https://upload.wikimedia.org/wikipedia/commons/3/39/Microsoft_timeline_of_operating_systems_2.png}
\caption{Arbre généalogique des systèmes Windows}
\end{figure}

\begin{itemize}
\item
  \textbf{Systèmes d'exploitation propriétaires :} Windows et
  partiellement MacOS (une partie UNIX est sous licence BSD)
\item
  \textbf{Systèmes d'exploitation libres :} de nombreuses distributions
  basées sur le noyau Linux, système dérivé d'UNIX : Debian, Ubuntu,
  Fedora, ArchLinux \ldots..
\end{itemize}

\begin{exercice}{}

\begin{enumerate}
\def\labelenumi{\arabic{enumi}.}
\tightlist
\item
  Lequel de ces systèmes d'exploitation est sous licence libre ?
\end{enumerate}

\textbf{\emph{Réponses}}

\textbf{A)} Android \textbf{B)} Linux \textbf{C)} Windows \textbf{D)}
MacOS

\begin{enumerate}
\def\labelenumi{\arabic{enumi}.}
\setcounter{enumi}{1}
\tightlist
\item
  Une et une seule de ces affirmations est \textbf{fausse}. Laquelle ?
\end{enumerate}

\textbf{\emph{Réponses}}

\textbf{A)} Un système d'exploitation libre est la plupart du temps
gratuit

\textbf{B)} Je peux contribuer à un système d'exploitation libre

\textbf{C)} Il est interdit d'étudier un système d'exploitation
propriétaire

\textbf{D)} Un système d'exploitation propriétaire est plus sécurisé

\end{exercice}

\end{document}
